% Cover letter for the MLST submission -- editor-facing, IDENTITY-BEARING.
% This file and its PDF must stay on the mlst-review strip list.
% Build: pdflatex cover_letter   (or `make cover` in paper/)
\documentclass[11pt]{article}
\usepackage{amsmath,amssymb}
\usepackage[hidelinks]{hyperref}
% ~1.1in margins without geometry.sty (not installed in this MiKTeX)
\setlength{\oddsidemargin}{0.1in}
\setlength{\textwidth}{6.3in}
\setlength{\topmargin}{-0.4in}
\setlength{\textheight}{9in}
% parskip-style paragraphs without parskip.sty
\setlength{\parindent}{0pt}
\setlength{\parskip}{1.8ex}
\pagestyle{empty}

\begin{document}

\begin{flushright}
Joshua Paul Brehm\\
Independent Researcher\\
\href{mailto:joshua.paul.brehm@gmail.com}{joshua.paul.brehm@gmail.com}\\[2ex]
July 30, 2026
\end{flushright}

\noindent The Editors\\
\emph{Machine Learning: Science and Technology}

\vspace{1ex}
\noindent\textbf{Re: Research-paper submission --- ``Heterogeneous Bipartite
Graph Neural Networks for Lattice Quantum Field Theory''}

\vspace{1ex}
Dear Editors,

Please consider the enclosed manuscript for publication in \emph{Machine
Learning: Science and Technology} as a research paper.

The manuscript's central contribution is architectural: a heterogeneous
bipartite graph representation for lattice field theories that separates
spacetime geometry from field content into distinct node types coupled by
typed edges --- a design inspired by the geometry--matter distinction in
continuum quantum field theory. Because the contribution is a representation
advance motivated by, and validated on, a physics problem, we believe MLST
--- at the interface of machine-learning methodology and the physical
sciences --- is its natural venue.

Four results support the design. First, we identify a structural failure
mode of the standard homogeneous-graph representation: message passing
destroys the field signal unless rescued by ad hoc skip connections, a
failure the bipartite design eliminates by construction. Second, the model
achieves near-perfect Euclidean-action prediction on two-dimensional scalar
$\varphi^4$ theory ($r > 0.9999$), against convolutional,
multilayer-perceptron, and homogeneous-GNN baselines. Third, a single model
trained at $16 \times 16$ transfers to lattices from $8 \times 8$ to
$64 \times 64$ with $r = 1.0000$ and no retraining. Fourth, the training
ensembles are validated by a finite-size-scaling study using a
Wolff/Brower--Tamayo cluster sampler, reproducing the exact Ising exponents
($\gamma/\nu = 1.732(8)$, $\nu = 1.04(8)$) once critical slowing down is
controlled --- itself an illustration of the sampling bottleneck that
motivates learned surrogates. The architecture accommodates gauge and
fermion fields at the same lattice sites without structural change, giving
a direct path toward lattice QCD.

We have selected \textbf{double-anonymous review}, and the manuscript is
anonymized per your checklist. The following declarations therefore appear
here rather than in the manuscript, to be restored to the Acknowledgements
on acceptance:

\begin{itemize}
  \item \textbf{Author:} Joshua Paul Brehm (independent researcher,
    unaffiliated), sole author, who takes full responsibility for the work.
  \item \textbf{Funding:} this work received no funding from any agency in
    the public, commercial, or not-for-profit sectors.
  \item \textbf{Conflicts of interest:} the author declares none.
  \item \textbf{Generative AI:} used in preparing the work and disclosed,
    with models, versions, and roles, in the manuscript's Acknowledgements
    per IOP policy (the disclosure is identity-neutral and retained in the
    review copy).
\end{itemize}

For reviewers, the complete code and analysis-level data --- sufficient to
regenerate every figure, table, and quoted number --- are available through
an anonymized repository mirror linked in the manuscript's Data availability
section; the public repository will be linked on acceptance.

The manuscript is original, has not been published previously, is not under
consideration elsewhere, and has not been posted to any preprint server.

Finally, as an unaffiliated researcher with no institutional or grant
support, I would be grateful to be considered for a waiver or discount of
the article publication charge, and I am happy to provide any documentation
required.

Thank you for your consideration.

\vspace{2ex}
Sincerely,

Joshua Paul Brehm

\end{document}
